% Petite bibliographie : les DOI des vingt contributions (générés depuis TOP20.json) et des leçons citées.
\chapter{Pour retrouver les travaux}
\markboth{Pour retrouver les travaux}{Pour retrouver les travaux}

Chaque travail cité dans ce livre a un DOI : tapez-le dans un moteur de recherche, ou ajoutez-le après
\texttt{https://doi.org/}, pour arriver à la page du dépôt, où se trouvent l'article, les données et le code.

\section{Les vingt contributions}
{\small
\setlist[description]{font=\normalfont\sffamily\bfseries\footnotesize\color{okbleu},leftmargin=3.2em,labelwidth=2.8em,
  labelsep=0.4em,itemsep=2pt}
% generated by code/make_backmatter.py from TOP20_final.json and chapters/ORDER.json -- do not edit
\begin{description}
\item[n°\,1] Des distances cosmiques avec une marge d'erreur vérifiée par une machine. {\footnotesize \doi{10.5281/zenodo.23247730}}
\item[n°\,2] Une tension célèbre de la cosmologie, écrite comme une inégalité vérifiée. {\footnotesize \doi{10.5281/zenodo.23284020}}
\item[n°\,3] Le triangle, rangement le moins coûteux : des conséquences vérifiées et une simulation. {\footnotesize \doi{10.5281/zenodo.23293429}}
\item[n°\,4] La répulsion écrantée entre dans la famille du théorème, et une erreur est réparée. {\footnotesize \doi{10.5281/zenodo.23646105}}
\item[n°\,5] Refaire le calcul d'un autre, sans le déclarer juste. {\footnotesize \doi{10.5281/zenodo.23653827}}
\item[n°\,6] Une règle prouvée droite montre que notre propre esquisse était fausse. {\footnotesize \doi{10.5281/zenodo.23653846}}
\item[n°\,7] Quelle chance a une mutation d'envahir une population~? {\footnotesize \doi{10.5281/zenodo.23653850}}
\item[n°\,8] Quand le plan de table ne change rien à la sélection. {\footnotesize \doi{10.5281/zenodo.23660071}}
\item[n°\,9] Un manuel pour suivre un théorème jusqu'au laboratoire. {\footnotesize \doi{10.5281/zenodo.23295142}}
\item[n°\,10] Des paires électron--trou qui se rangent en triangles, sur le papier. {\footnotesize \doi{10.5281/zenodo.23646103}}
\item[n°\,11] La force dans un semi-conducteur atomiquement mince, en tas de sable. {\footnotesize \doi{10.5281/zenodo.23646107}}
\item[n°\,12] Tester avant de prouver : une idée naturelle qui ne tient pas. {\footnotesize \doi{10.5281/zenodo.23653833}}
\item[n°\,13] Parfait au repos, pas quand il frissonne. {\footnotesize \doi{10.5281/zenodo.23655392}}
\item[n°\,14] Quatre étages sur douze, et le prix du reste. {\footnotesize \doi{10.5281/zenodo.23656804}}
\item[n°\,15] Remonter la machine des autres dans son propre atelier. {\footnotesize \doi{10.5281/zenodo.23661381}}
\item[n°\,16] Une garantie générale, trois certificats précis. {\footnotesize \doi{10.5281/zenodo.23285631}}
\item[n°\,17] Le moteur des calculs, et son rapport d'inspection. {\footnotesize dépôt \texttt{rusty-SUNDIALS}}
\item[n°\,18] Une porte qui reste fermée quand on ne sait pas. {\footnotesize \doi{10.5281/zenodo.23119953}, \doi{10.5281/zenodo.23120120}}
\item[n°\,19] Le banc d'essai qui s'est inspecté lui-même. {\footnotesize dépôt \texttt{AutoevolveAI}}
\item[n°\,20] On accusait le thermostat de surchauffer. {\footnotesize \doi{10.5281/zenodo.23478809}, \doi{10.5281/zenodo.23480005}, \doi{10.5281/zenodo.23481378}, \doi{10.5281/zenodo.23482410}, \doi{10.5281/zenodo.23483052}, \doi{10.5281/zenodo.23483564}, \doi{10.5281/zenodo.23483643}, \doi{10.5281/zenodo.23532458}}
\end{description}

}

\section{Les leçons et les méthodes citées}
{\small
\begin{description}[font=\normalfont\sffamily\bfseries\footnotesize\color{okvermillon},leftmargin=0pt,itemsep=2pt]
\item[La méthode] « A laboratory method for AI-assisted science » (18 septembre 2026). \doi{10.5281/zenodo.22829767}
\item[RunuX, version 2] retrait des résultats jamais obtenus. \doi{10.5281/zenodo.22879667}
\item[Dual-scale, révision 5] la théorie confrontée aux données et retirée. \doi{10.5281/zenodo.22837842}
\item[Navier--Stokes, version 2] l'interprétation retirée. \doi{10.5281/zenodo.22773907}
\item[RunuX, audit] affirmations contre preuves. \doi{10.5281/zenodo.22985926}
\item[SQIsign, édition 2] trois affirmations retirées. \doi{10.5281/zenodo.23532601}
\item[Thermodynamique stochastique, édition 2] la mention « relu par des pairs » retirée. \doi{10.5281/zenodo.23525655}
\item[Hopfield--Ninio, édition 3] le mécanisme de l'article réfuté. \doi{10.5281/zenodo.23526783}
\item[QU-1] la réfutation numérique. \doi{10.5281/zenodo.23653833}
\end{description}
}
